\documentclass[12pt]{article}

% Modelo de relatório do Trabalho 01, baseado no template de conferências da
% Sociedade Brasileira de Computação (SBC), versão atualizada em 2017.
\usepackage{sbc-template}
\usepackage{graphicx,url}
\usepackage[utf8]{inputenc}
\usepackage[T1]{fontenc}
\usepackage{amsmath,amssymb}
\usepackage{booktabs}
\usepackage{tabularx}
\usepackage{multirow}
\usepackage{float}
\usepackage[hidelinks]{hyperref}

\sloppy
\renewcommand{\tablename}{Tabela}
\renewcommand{\figurename}{Figura}
\renewcommand{\refname}{Referências}

\title{Trabalho 01 - Busca Não Informada, Busca Informada e Busca Local}

% Substitua pelos nomes da dupla. Se o trabalho for individual, remova o segundo nome.
\author{Lorena Ávila de Souza Silva\inst{1}}

\address{Curso de Engenharia de Computação - CEFET-MG, campus Divinópolis\\
Disciplina de Inteligência Artificial - 2026/2
\email{lorena.silva@aluno.cefetmg.br}}

\begin{document}

\maketitle

\begin{abstract}
This report presents the implementation and experimental analysis of uninformed, informed, and local search algorithms developed for Assignment 01. The experiments compare BFS, DFS, Uniform-Cost Search, Greedy Best-First Search, and A* in a weighted grid, and Hill Climbing and Random-Restart Hill Climbing in the 8-Queens problem. The results demonstrate the trade-offs between path length, computational cost, memory overhead, and heuristic efficiency.
\end{abstract}

\begin{resumo}
Este relatório apresenta a implementação e a análise experimental dos algoritmos de busca desenvolvidos no Trabalho 01. Os experimentos comparam BFS, DFS, Busca de Custo Uniforme, Busca Gulosa e A* em um ambiente ponderado, além de Hill Climbing e Random-Restart Hill Climbing no problema das 8 rainhas. Os resultados evidenciam as trocas entre comprimento de caminho, custo computacional, uso de memória e a eficácia das heurísticas avaliadas.
\end{resumo}

\section{Introdução}
O presente trabalho aborda a modelagem e a resolução de problemas clássicos de Inteligência Artificial divididos em duas frentes: a busca em ambientes ponderados em grade (caminhos ótimos e subótimos) e a otimização local aplicada ao problema das 8 rainhas. O objetivo principal é comparar o comportamento prático de diferentes estratégias de busca (não informada, informada e local) frente a métricas rigorosas de desempenho, como número de passos, custo acumulado, nós gerados, nós expandidos, picos de memória e tempo de execução.

Ao final desta seção, indica-se que a análise busca responder a questões fundamentais sobre a divergência entre menor número de passos e menor custo real, o impacto da escolha da heurística (Manhattan vs. Euclidiana) na eficiência do A* e da Busca Gulosa, e as limitações de ótimos locais enfrentadas pelo Hill Climbing puro em contraste com o sucesso do Random Restart.

\section{Representação e implementação}
\subsection{Parte 1 - Busca não informada e informada}
A representação do ambiente consiste em uma grade ponderada onde cada célula possui um custo associado de transição. As buscas (BFS, DFS, UCS, Gulosa e A*) utilizam estruturas de dados nativas eficientes: `deque` para a fila da BFS, listas operando como pilhas para a DFS, e filas de prioridade (`heapq`) para UCS, Gulosa e A*, garantindo desempates determinísticos baseados em contadores de inserção e critérios específicos de avaliação. 

O rastreamento de estados descobertos foi implementado no momento da inserção na fronteira (para BFS e DFS) e por meio de dicionários de custos ótimos (`cost_so_far` / `g_score`) para UCS e A*, permitindo o descarte eficiente de entradas obsoletas antes do teste de objetivo. A reconstrução do caminho é feita retroativamente através de um dicionário de pais (`parent`), isolando a lista de posições (`path`) e de ações tomadas (`actions`), culminando no cálculo do custo total via soma dos custos de passo a partir de `path[1:]`.

\subsection{Heurísticas}
As heurísticas utilizadas foram as distâncias de Manhattan e Euclidiana em relação ao objetivo. Ambas são adequadas ao modelo de movimento em quatro direções ortogonais (Norte, Sul, Leste, Oeste) com custos unitários mínimos. A admissibilidade e consistência de ambas decorrem de sua formulação geométrica sobre grades, onde nunca superestimam o custo real restante para alcançar o objetivo.

\subsection{Problema das 8 rainhas}
O problema foi modelado representando uma configuração como um vetor de 8 inteiros, onde cada índice indica a coluna e o valor armazena a linha da rainha correspondente. A função de avaliação $h(s)$ calcula o número total de pares de rainhas em conflito (mesma linha ou mesma diagonal). A vizinhança de tamanho 56 é gerada explorando todas as 8 colunas e as 7 posições alternativas de linha para cada uma. 

O Hill Climbing escolhe de forma estritamente gulosa o primeiro vizinho que reduza o valor de $h$, recusando movimentos laterais ou plateaus e parando quando nenhum vizinho melhora o estado atual. O Random Restart reutiliza o gerador pseudo-aleatório (`random.Random` com sementes de 0 a 99) para reinicializar o tabuleiro até 50 vezes em caso de falha.

\section{Metodologia experimental}
Os experimentos foram executados em um ambiente Linux (Ubuntu) utilizando Python 3, sem dependências externas além das bibliotecas padrão. O mapa público utilizado na Parte 1 foi o `data/mapa_teste.txt`. As métricas de tempo foram obtidas isolando a leitura do arquivo e calculando a mediana de 30 execuções independentes (`time.perf_counter()`).

Para a Parte 2, executou-se o Hill Climbing puro e o Random Restart sobre todas as 100 sementes inteiras de 0 a 99, respeitando o limite máximo de 50 reinícios. O código-fonte completo e reprodutível encontra-se no repositório Git oficial.

\noindent\textbf{Repositório:} \url{https://github.com/seu-usuario/ia-trabalho01} \\
\textbf{Tag da entrega final:} \texttt{entrega-final}.

\section{Resultados}
\subsection{Parte 1 - Busca não informada e informada}
A Tabela~\ref{tab:grid} consolida os resultados obtidos no mapa de teste público.

\begin{table}[H]
\centering
\caption{Resultados dos algoritmos de busca no mapa de teste.}
\label{tab:grid}
\resizebox{\textwidth}{!}{%
\begin{tabular}{llrrrrrrr}
\toprule
Algoritmo & Heurística & Passos & Custo & Gerados & Expandidos & Pico fronteira & Memória & Tempo (ms) \\
\midrule
BFS     & -          & 15 & 40 & 66 & 59 & 9  & 66 & 0.123 \\
DFS     & -          & 15 & 40 & 21 & 15 & 6  & 21 & 0.042 \\
UCS     & -          & 29 & 29 & 77 & 74 & 6  & 77 & 0.189 \\
Gulosa  & Manhattan  & 15 & 40 & 21 & 15 & 6  & 21 & 0.040 \\
Gulosa  & Euclidiana & 15 & 40 & 21 & 15 & 6  & 21 & 0.042 \\
A*      & Manhattan  & 29 & 29 & 60 & 43 & 15 & 58 & 0.132 \\
A*      & Euclidiana & 29 & 29 & 72 & 60 & 12 & 72 & 0.231 \\
\bottomrule
\end{tabular}}
\end{table}

\subsection{Parte 2 - Busca local e 8 rainhas}
A Tabela~\ref{tab:queens} resume os experimentos realizados com as 100 sementes (0 a 99).

\begin{table}[H]
\centering
\caption{Resultados dos métodos de busca local no problema das 8 rainhas.}
\label{tab:queens}
\begin{tabularx}{\textwidth}{Xrr}
\toprule
Métrica & Hill Climbing & Random Restart \\
\midrule
Execuções & 100 & 100 \\
- Soluções encontradas & 12 & 100 \\
- Taxa de sucesso (\%) & 12\% & 100\% \\
- Média de iterações & 3.28 & 24.72 \\
- Mediana de iterações & 3.0 & 17.5 \\
- Média de estados avaliados & 240.68 & 1826.64 \\
- Média de reinícios & N/A & 6.76 \\
- Máximo de reinícios & N/A & 42 \\
\bottomrule
\end{tabularx}
\end{table}

Um exemplo típico de travamento do Hill Climbing puro ocorreu na seed 0, onde o algoritmo estagnou em um ótimo local com $h = 2$ e tabuleiro final `[6, 3, 5, 2, 1, 6, 4, 7]`, onde nenhum dos 56 estados vizinhos era estritamente superior (menor $h$).

\section{Discussão}
A análise conjunta dos resultados permite responder diretamente às questões centrais propostas:
\begin{itemize}
    \item \textbf{Menor número de ações vs. menor custo:} A BFS encontrou um caminho com menor número de passos (15 ações), mas custo elevado (40), enquanto o UCS encontrou um caminho mais longo em passos (29) mas com o menor custo ótimo (29), provando que caminhos curtos podem atravessar terrenos de alto custo.
    \item \textbf{Buscas não informadas:} A BFS garante menor caminho em grafos não ponderados, mas falha em otimizar custos. A UCS expandiu mais nós (74) para garantir a optimalidade do custo. A DFS encontrou solução por sorte estrutural, mas sem garantias.
    \item \textbf{Busca Gulosa:} Mostrou-se extremamente rápida e econômica em expansões (15 nós), mas é subótima (custo 40), pois prioriza cegamente a proximidade heurística do objetivo sem acumular o custo do trajeto.
    \item \textbf{Comparação UCS e A*:} O A* com Manhattan reduziu as expansões de 74 (UCS) para 43, mantendo o mesmo custo ótimo de 29, demonstrando a eficácia da poda guiada por heurística admissível.
    \item \textbf{Manhattan vs. Euclidiana:} No A*, Manhattan expandiu 43 nós contra 60 da Euclidiana. Como o movimento é em grade retilínea (4 direções), Manhattan fornece uma estimativa mais próxima e informativa do custo real do que a linha reta da Euclidiana.
    \item \textbf{Hill Climbing e Random Restart:} O HC puro apresentou baixa taxa de sucesso (12\%) devido a ótimos locais. O mecanismo de Random Restart elevou o sucesso para 100% ao permitir saltar de armadilhas topográficas utilizando novas sementes aleatórias.
\end{itemize}

\section{Conclusão}
O trabalho permitiu consolidar a compreensão prática entre buscas cegas, informadas e locais. Ficou evidente que o A* com heurística de Manhattan atinge o melhor equilíbrio de optimalidade e eficiência em ambientes ponderados, enquanto o Random Restart é indispensável para garantir a completude prática de métodos gulosos de subida de encosta em espaços complexos.

\section*{Contribuições dos autores}
Todas as etapas de implementação, testes, execução de experimentos e redação deste relatório foram realizadas de forma individual pela autora.

\section*{Declaração de uso de IA}
Foram utilizadas ferramentas de IA de forma assistiva para verificação de sintaxe de testes unitários e esclarecimento de pequenos erros de formatação de código, conforme detalhado no arquivo \texttt{IA\_USAGE.md}.

\begin{thebibliography}{99}

\bibitem[Stuart and Norvig 2021]{russell2021}
Russell, S. and Norvig, P. (2021). \textit{Artificial Intelligence: A Modern Approach}. Pearson, 4th edition.

\end{thebibliography}

\end{document}